\section{Distribusi geometrik}
\label{geomDist}

Berapa kali, menurut perkiraan kita, sebuah koin perlu dilempar sampai
muncul sisi \resp{kepala}? Atau berapa kali sebuah dadu perlu dilempar
sampai kita memperoleh \resp{1}? Pertanyaan-pertanyaan ini dapat dijawab
dengan distribusi geometrik. Pertama-tama, kita memformalkan setiap
percobaan--seperti satu lemparan koin atau dadu--dengan distribusi
Bernoulli. Kemudian, kita menggabungkannya dengan perangkat peluang dari
Bab~\ref{probability} untuk membangun distribusi geometrik.

\subsection{Distribusi Bernoulli}
\label{bernoulli}

\newcommand{\insureSprob}{0.7}
\newcommand{\insureSperc}{70\%}
\newcommand{\insureFprob}{0.3}
\newcommand{\insureFperc}{30\%}
\newcommand{\insureDistA}{0.7}
\newcommand{\insureDistB}{0.21}
\newcommand{\insureDistC}{0.063}
\newcommand{\insureDistD}{0.019}
\newcommand{\insureDistE}{0.006}
\newcommand{\insureCDistA}{0.7}
\newcommand{\insureCDistB}{0.91}
\newcommand{\insureCDistC}{0.973}
\newcommand{\insureCDistCComplement}{0.027}
\newcommand{\insureCDistD}{0.992}
\newcommand{\insureCDistE}{0.998}
\newcommand{\insureGeomMean}{1.43}

\index{distribution!Bernoulli|(}

Banyak program asuransi kesehatan di Amerika Serikat menerapkan
risiko sendiri (\emph{deductible}). Dalam skema ini, tertanggung menanggung
biaya hingga batas risiko sendiri; biaya di atas batas itu kemudian dibagi
antara tertanggung dan perusahaan asuransi selama sisa tahun berjalan.

Misalkan sebuah perusahaan asuransi kesehatan mendapati bahwa
\insureSperc{} orang yang diasuransikannya tetap berada di bawah batas
risiko sendiri pada suatu tahun. Setiap orang tersebut dapat dipandang sebagai
satu \termsub{percobaan}{trial}. Kita memberi label
\termsub{sukses}{success} kepada seseorang jika biaya layanan kesehatannya
tidak melampaui batas risiko sendiri. Kita memberi label
\termsub{gagal}{failure} jika biaya orang tersebut melampaui batas risiko
sendiri pada tahun itu. Karena 70\% orang tidak mencapai batas risiko
sendiri, kita menyatakan \termsub{peluang sukses}{probability of a success}
sebagai $p = \insureSprob{}$. Peluang gagal kadang-kadang dinyatakan dengan
$q = 1 - p$, yang bernilai \insureFprob{} pada contoh asuransi ini.

Jika sebuah percobaan hanya memiliki dua hasil yang mungkin, yang sering
diberi label \resp{sukses} atau \resp{gagal}, percobaan itu disebut
\termsub{variabel acak Bernoulli}{distribution!Bernoulli}.
Kita memilih untuk memberi label ``sukses'' kepada orang yang tidak
mencapai batas risiko sendiri dan ``gagal'' kepada semua orang lainnya.
Namun, kedua label tersebut dapat saja kita pertukarkan. Kerangka
matematis yang akan kita bangun tidak bergantung pada hasil mana yang
diberi label sukses dan mana yang diberi label gagal, asalkan kita
konsisten.

Variabel acak Bernoulli sering dinyatakan dengan \resp{1} untuk sukses
dan \resp{0} untuk gagal. Selain memudahkan pemasukan data, pelabelan ini
juga praktis secara matematis. Misalkan kita mengamati sepuluh percobaan:
\begin{center}
\resp{1} \resp{1} \resp{1} \resp{0} \resp{1} \resp{0} \resp{0} \resp{1} \resp{1} \resp{0}
\end{center}
Maka \termsub{proporsi sampel}{sample proportion}, $\hat{p}$, adalah
rata-rata sampel pengamatan-pengamatan tersebut:
\begin{align*}
\hat{p} = \frac{\text{\# sukses}}{\text{\# percobaan}}
    = \frac{1+1+1+0+1+0+0+1+1+0}{10} = 0.6
\end{align*}%
Kajian matematis terhadap variabel acak Bernoulli ini masih dapat
diperluas.
%\Comment{Maybe the next footnote should instead be an EOCE?}
Karena \resp{0} dan \resp{1} merupakan hasil numerik, kita dapat
mendefinisikan {rata-rata} dan {simpangan baku} variabel acak Bernoulli.
(Lihat Latihan~\ref{bernoulli_mean_derivation}
dan~\ref{bernoulli_sd_derivation}.)

\begin{onebox}{Variabel acak Bernoulli}
%  A Bernoulli random variable has exactly two possible
%  outcomes, often labeled \resp{1} for the ``success''
%  outcome and \resp{0} for the ``failure'' outcome.\vspace{3mm}
  Jika $X$ adalah variabel acak yang bernilai 1 dengan
  peluang sukses $p$ dan bernilai 0 dengan peluang $1-p$,
  maka $X$ adalah variabel acak Bernoulli dengan rata-rata
  dan simpangan baku
  \begin{align*}
  \mu &= p
      &\sigma&= \sqrt{p(1-p)}
  \end{align*}
\end{onebox}

Secara umum, variabel acak Bernoulli berguna jika dipandang sebagai proses
acak yang hanya memiliki dua hasil: sukses atau gagal. Selanjutnya, kita
membangun kerangka matematis menggunakan label numerik \resp{1} dan
\resp{0}, masing-masing untuk sukses dan gagal.

\index{distribution!Bernoulli|)}


\D{\newpage}

\subsection{Distribusi geometrik}

\index{distribution!geometric|(}

\termsub{Distribusi geometrik}{distribution!geometric} digunakan untuk
menggambarkan berapa banyak percobaan yang diperlukan sampai suatu sukses
diamati. Mari kita awali dengan sebuah contoh.

\begin{examplewrap}
\begin{nexample}{Misalkan kita bekerja di perusahaan asuransi dan
    memerlukan satu kasus orang yang tidak melampaui batas risiko sendiri
    sebagai studi kasus. Jika peluang seseorang tidak melampaui batas
    risiko sendiri adalah \insureSprob{} dan kita memilih orang secara acak,
    berapa peluang bahwa orang pertama tidak melampaui batas risiko sendiri,
    yakni menjadi sukses? Bagaimana dengan orang kedua? Orang ketiga?
    Bagaimana jika kita memilih $n - 1$ kasus sebelum menemukan sukses
    pertama, yakni sukses pertama ditemukan pada orang ke-$n$?
    (Jika sukses pertama ditemukan pada orang kelima, kita menyatakan
    $n=5$.)}
  \label{waitForDeductible}%
  Peluang bahwa kita berhenti setelah orang pertama sama dengan peluang
  orang pertama tidak mencapai batas risiko sendiri:~\insureSprob{}.
  Peluang bahwa orang kedua merupakan orang pertama yang tidak mencapai
  batas risiko sendiri adalah
  \begin{align*}
  &P(\text{orang kedua menjadi sukses pertama}) \\
  &\quad
    = P(\text{orang pertama gagal, orang kedua sukses})
    = (\insureFprob{})(\insureSprob{})
    = \insureDistB{}
  \end{align*}
  Demikian pula, peluang bahwa sukses pertama ditemukan pada kasus ketiga
  adalah $(\insureFprob{})(\insureFprob{})(\insureSprob{})
    = \insureDistC$.

  Jika sukses pertama ditemukan pada orang ke-$n$, terdapat $n-1$
  kegagalan yang diikuti 1 sukses. Urutan ini bersesuaian dengan peluang
  $(\insureFprob{})^{n-1}(\insureSprob{})$.
  Nilai ini sama dengan
  $(1-\insureSprob{})^{n-1}(\insureSprob{})$.
\end{nexample}
\end{examplewrap}

Contoh~\ref{waitForDeductible} mengilustrasikan
\termsub{distribusi geometrik}{distribution!geometric}, yang menggambarkan
waktu tunggu hingga sukses untuk variabel acak Bernoulli yang
\termsub{saling independen dan berdistribusi identik (iid)}{independent and identically distributed (iid)}.
Dalam kasus ini, aspek \emph{independensi} hanya berarti
bahwa setiap orang dalam contoh tidak saling memengaruhi, sedangkan
\emph{identik} berarti bahwa masing-masing memiliki peluang sukses yang
sama.

Distribusi geometrik dari Contoh~\ref{waitForDeductible} ditampilkan pada
Gambar~\ref{geometricDist70}. Secara umum, peluang-peluang pada distribusi
geometrik menurun dengan cepat \termsub{secara eksponensial}{exponentially}.

\begin{figure}[h]
  \centering
  \Figure[Ditampilkan distribusi peluang ``Jumlah Percobaan hingga Sukses untuk p = 0.7'' dalam bentuk diagram batang. Nilai yang ditampilkan adalah 1, 2, 3, 4, 5, 6, 7, dan 8. Peluangnya masing-masing sekitar 0.7, 0.21, 0.07, 0.02, dan 0.01; untuk nilai 6, 7, dan 8, tinggi batang hampir tidak dapat dibedakan.]{0.8}{geometricDist70}
  \caption{Distribusi geometrik ketika peluang sukses adalah
      $p = \insureSprob{}$.}
  \label{geometricDist70}
\end{figure}

Teks ini tidak menurunkan rumus rata-rata (nilai harapan) banyaknya
percobaan yang diperlukan untuk menemukan sukses pertama, ataupun rumus
simpangan baku dan varians distribusi ini. Namun, kita menyajikan rumus
umum untuk ketiganya.

\begin{onebox}{Distribusi geometrik}
  \index{distribution!geometric|textbf}%
  Jika peluang sukses dalam satu percobaan adalah $p$ dan peluang gagal
  adalah $1-p$, peluang bahwa sukses pertama ditemukan pada percobaan
  ke-$n$ diberikan oleh\vspace{-1.5mm}
  \begin{align*}
  (1-p)^{n-1}p
  \end{align*}
  Rata-rata (yakni nilai harapan), varians, dan simpangan baku waktu
  tunggu ini diberikan oleh
  \begin{align*}
  \mu &= \frac{1}{p}
      &\sigma^2 &=\frac{1-p}{p^2}
      &\sigma &= \sqrt{\frac{1-p}{p^2}}
  \end{align*}
\end{onebox}

Bukan kebetulan bahwa kita menggunakan simbol $\mu$ untuk rata-rata
sekaligus nilai harapan. Rata-rata dan nilai harapan adalah hal yang sama.

Dalam distribusi geometrik, rata-rata banyaknya percobaan yang diperlukan
untuk memperoleh sukses adalah $1/p$. Hasil matematis ini sesuai dengan
intuisi. Jika peluang sukses tinggi (misalnya 0.8), biasanya kita tidak
menunggu lama untuk memperoleh sukses: rata-rata $1/0.8 = 1.25$ percobaan.
Jika peluang sukses rendah (misalnya 0.1), kita memperkirakan perlu
mengamati banyak percobaan sebelum melihat sukses: $1/0.1 = 10$ percobaan.

\begin{exercisewrap}
\begin{nexercise}
Peluang bahwa suatu kasus tidak melampaui batas risiko sendiri adalah
\insureSprob{}. Jika kita memeriksa kasus satu demi satu sampai menemukan
seseorang yang tidak mencapai batas risiko sendiri, berapa banyak kasus
yang diperkirakan perlu kita periksa?\footnotemark{}
\end{nexercise}
\end{exercisewrap}
\footnotetext{Kita memperkirakan perlu memeriksa sekitar
    $1 / \insureSprob{} \approx \insureGeomMean{}$
    orang untuk menemukan sukses pertama.}

\begin{examplewrap}
\begin{nexample}{Berapa peluang bahwa kita menemukan sukses pertama
    di antara tiga kasus pertama?}
  \label{insureFirstSuccessInLT4}%
  Sukses pertama dapat ditemukan pada kasus pertama ($n=1$), kedua
  ($n=2$), atau ketiga ($n=3$). Ketiganya merupakan hasil yang saling
  lepas. Karena orang-orang dalam sampel dipilih secara acak dari populasi
  besar, pemilihan mereka saling independen. Kita menghitung peluang
  setiap kasus lalu menjumlahkan hasil-hasil yang terpisah tersebut:
  \begin{align*}
  &P(n=1, 2, \text{ atau }3) \\
    & \quad = P(n=1)+P(n=2)+P(n=3) \\
    & \quad = (\insureFprob{})^{1-1}(\insureSprob{})
        + (\insureFprob{})^{2-1}(\insureSprob{})
        + (\insureFprob{})^{3-1}(\insureSprob{}) \\
    & \quad = \insureCDistC{}
  \end{align*}
  Peluang bahwa kita menemukan kasus yang sukses dalam tiga kasus pertama
  adalah \insureCDistC{}.
\end{nexample}
\end{examplewrap}

\begin{exercisewrap}
\begin{nexercise}
Tentukan cara yang lebih cerdik untuk menyelesaikan
Contoh~\ref{insureFirstSuccessInLT4}. Tunjukkan bahwa Anda memperoleh hasil
yang sama.\footnotemark{}
\end{nexercise}
\end{exercisewrap}
\footnotetext{Pertama, cari peluang komplemennya:
  $P($tidak ada sukses dalam 3~percobaan pertama$)
      = \insureFprob{}^3 = \insureCDistCComplement{}$.
  Selanjutnya, hitung satu dikurangi peluang tersebut:
  $1 - P($tidak ada sukses dalam 3 percobaan$)
      = 1 - \insureCDistCComplement{}
      = \insureCDistC{}$.}

\D{\newpage}

\begin{examplewrap}
\begin{nexample}{Misalkan sebuah perusahaan asuransi mobil telah
    menentukan bahwa 88\% pengemudinya tidak melampaui batas risiko sendiri
    pada suatu tahun. Jika seorang pegawai perusahaan memilih berkas
    pengemudi secara acak sampai menemukan satu pengemudi yang tidak
    melampaui batas risiko sendiri, berapa banyak berkas pengemudi yang
    diperkirakan perlu diperiksa pegawai tersebut? Berapa simpangan baku
    banyaknya berkas pengemudi yang harus dipilih?}
  \label{carInsure08DrawOne}%
  Dalam contoh ini, sukses kembali didefinisikan sebagai keadaan ketika
  seseorang tidak melampaui batas risiko sendiri pada polisnya, yang memiliki
  peluang $p = 0.88$. Banyaknya orang yang diperkirakan perlu diperiksa
  adalah $1 / p = 1 / 0.88 = 1.14$, dan simpangan bakunya adalah
  $\sqrt{(1-p)/p^2} = 0.39$.
\end{nexample}
\end{examplewrap}

\begin{exercisewrap}
\begin{nexercise}
Berdasarkan hasil dari Contoh~\ref{carInsure08DrawOne}, yaitu
$\mu = 1.14$ dan $\sigma = 0.39$, apakah model normal tepat digunakan
untuk mencari proporsi eksperimen yang berakhir dalam 3 percobaan atau
kurang?\footnotemark{}
\end{nexercise}
\end{exercisewrap}
\footnotetext{Tidak. Distribusi geometrik selalu menceng kanan dan tidak
  pernah dapat didekati dengan baik oleh model normal.}

Asumsi independensi sangat penting agar distribusi geometrik dapat
menggambarkan suatu skenario secara akurat. Secara matematis, ketika
membangun peluang sukses pada percobaan ke-$n$, kita harus menggunakan
Aturan Perkalian untuk proses yang saling independen. Tidak mudah
menggeneralisasi model geometrik untuk percobaan-percobaan yang saling
dependen.

\index{distribution!geometric|)}


{\input{ch_distributions/TeX/geometric_distribution.tex}}
